\documentclass[11pt]{article}
\usepackage{kotex,amsmath,amssymb,booktabs}
\usepackage[margin=1in]{geometry}
\usepackage[hidelinks]{hyperref}
\newcommand{\norm}[1]{\left\lVert #1\right\rVert}
\title{Causal Allocation Editing: 실제 MEMIT write의 공동 최적화와 층별 배분}
\author{Method 설계 및 구현 계약}
\date{2026년 10월 6일}
\begin{document}
\maketitle

\part*{현재 설계: Causal Allocation Editing}
\paragraph{범위와 실행 상태.}
Native MEMIT형 mean-key ridge writer와 기존 history를 유지하면서,
모든 후보층의 local 요청량과 배분을 하나의 실제 후보 모델에서 최적화한다.
사용자가 선택한 all-token actual forward와 full causal gradient를 유지한다.
2026년 10월 6일 사용자는 실험명을 새로 정하여 GH에게 전달하고
SH4가 실행하도록 지시했다. 실험명은 \textbf{Causal Allocation Editing},
task/job/result 식별자는 \texttt{causal-allocation-editing}이며 프로젝트 접두사와
버전 접미사를 넣지 않는다. 내부 문서 revision과 과거 경로는 재현 이력으로만 남긴다.
Production 구현, 실제 모델 검증과 GPU 제출은 실행 담당자의 후속 작업이다.
GH $\to$ SH4/server4의 단일 MEMIT arm, 2,000편집 범위는
\texttt{../2026-10-06-jlz-causal-joint-writer/experiment-2k.json}에 보존한다.
위 \texttt{2026-10-06-jlz-causal-joint-writer/} 디렉터리의
method/implementation/validation JSON이 구현 계약이다.
이전 R1 계약은 그 디렉터리 아래 \texttt{history/r1-contracts.json}에 보존했다.
아래 CD-Q/CD-C, V14 및 V13은 역사 기록이며 현재 규칙에 섞지 않는다.

\section{검증하려는 claim과 주 학습 경로}
교정하는 문제는 두 가지다.
첫째, local 목표의 효과를 평가하는 상태와 실제 순차 write 상태를 같게 한다.
둘째, 하층 write가 상층 key와 writer 비용을 바꾸는 경로까지 미분하여
층별 요청량을 공동으로 결정한다. 모든 설정된 편집층은 항상 후보이며,
특정 층으로의 집중과 정확한 0 배분 모두 허용한다.

Native rewrite/prefix/KL 문장, lookup, target, tokenization,
NLL reduction과 readout, full-vocabulary KL(current$\Vert$entry),
rewrite-only native key 평균 및 history 의미는 보존한다.
별도 reference dataset, history replay, 추가 paraphrase/generation 문장,
subject 누수 목적, exact writer, 역실현 증폭과 absolute-z tracking은 넣지 않는다.
단, native compute-z의 subject delta hook 대신
\textbf{배포될 후보 weight가 모든 token에 작용한 forward}에서 loss를 계산한다.
데이터와 loss 식을 보존한다는 사실이 native 전체 학습 연산자와 같다는 뜻은 아니다.

\paragraph{Routing과 frozen-upper의 위치.}
주 경로는 수정하지 않은 목적의 full gradient다.
DOW-KE의 NLL pre-subject direct-parameter gradient routing은
input-key gradient와 KL을 유지하는 별도 backward 규칙이다.
우리의 candidate-dependent $P(K)$에 옮기면 stop-gradient 위치에 따라
서로 다른 gradient가 되므로 이번 main에는 넣지 않는다.
이는 \href{https://arxiv.org/html/2608.16932v1}{DOW-KE 원문}의
방법과 우리 full-gradient 계약을 구분한 결정이다.
상층 key 고정도 현재의 두 교정을 근사로 바꾸므로 main에서 제외한다.
과거 two-arm의 PS 약 67.9\%는 여러 설계가 함께 달랐던 결과이며
all-token forward만의 실패를 입증하지 않는다.

\section{변수, 층별 native norm과 상한}
Batch entry의 $W^e,H^e,C_0$, teacher 및 canonical anchor를 고정한다.
각 층 $l$, 요청 $r$에서 $a_{lr}>0$를 native subject hidden norm이라 두고
\[
 R_{lr}=a_{lr}u_{lr},\qquad
 \norm{u_{lr}}_2\le c_{lr},\qquad c_{lr}=0.75
\]
로 정의한다. 일반 모델에서는 adapter의 native profile 값을 쓴다.
$R$은 자기 write 층의 requested residual이며 실현을 보장하는 양이 아니다.
모든 $u$는 0에서 시작한다.
\[
 G(u)=\sum_{l,r}\lambda_n\frac{\norm{R_{lr}}_2}{a_{lr}^2}
     =\sum_{l,r}\beta_{lr}\norm{u_{lr}}_2,\qquad
 \beta_{lr}=\frac{\lambda_n}{a_{lr}},\quad \lambda_n=0.5 .
\]
기존 공통 $a_\star$와 요청별 공유 $\sum_l\norm{u_{lr}}\le.75$는 제거한다.
총 요청량의 허용 영역은 커진다. 위 제약은 층별 상한일 뿐,
총 실현량이나 의미적 강도를 native와 같게 보장하지 않는다.
Task, $G$, $Q$, 상한과 유한한 solver 경로가 함께 크기와 배분을 결정한다.
같은 $0.5$라도 다층 prox의 작용이 단층 native Adam과 같다고 주장하지 않는다.

\section{후보 하나의 causal MEMIT write}
각 후보는 같은 $W^e/H^e$에서 시작한다.
하층의 후보 weight가 적용된 실제 상태에서 층 $l$의 native subject key를 읽고,
rewrite-only nested mean으로 논리 batch 전체의 $K_l(u_{<l})$를 만든다.
KL 문장은 loss에 쓰되 writer key 평균에는 넣지 않는다.
\[
 A_l=\lambda_{C,l}C_{0,l}+H_l^e,\qquad
 P_l=(A_l+K_lK_l^\top)^{-1}K_l,\qquad U_l=R_lP_l^\top ,
\]
\[
 W_l^{\rm cand}=W_l^e+\operatorname{cast}_{32}(U_l).
\]
현재 profile은 $\lambda_C=15000$, FP32 모델, FP64 geometry다.
역행렬은 명시적으로 만들지 않고 native-equivalent solve를 사용한다.
Own key는 own write 이전 입력이므로 순환 정의가 없다.
같은 weight로 모든 token을 처리한 뒤 상층 key를 다시 구한다.
물리 microbatch마다 별도 writer를 풀지 않는다.

Local 요청 목표는 개념적으로
$z_{lr}^{\rm req}=h_{lr}^{\rm actual,pre}+R_{lr}$다.
저장한 $R$을 그대로 RHS로 사용하며, 다른 가상 상태의 $z-h$를 빼어
계획 밖 상속 보정량을 만들지 않는다.
손실과 commit이 관측하는 effective action은
\[
 Y_{lrc}=U_l^{\rm eff}k_{lrc}^{\rm actual},\qquad
 U_l^{\rm eff}=\operatorname{double}(W_l^{\rm cand})
             -\operatorname{double}(W_l^e)
\]
다. $R=0$인 열도 solve에 남고, $R_l=0$인 층의 gradient도 남아 재진입할 수 있다.

\section{단일 목적과 완전한 미분}
\[
 \min_{\norm{u_{lr}}\le c_{lr}} J(u)=F(u)+G(u),
\]
\[
 F(u)=\sum_r\left[\mathrm{NLL}_r(W^{\rm cand}(u))
 +\lambda_{\rm KL}\mathrm{KL}_r(\mathrm{current}(u)\Vert\mathrm{entry})\right]
 +\lambda_Q\sum_lQ_l(u),\qquad \lambda_{\rm KL}=0.0625 ,
\]
\[
 Q_l=\operatorname{tr}(U_lA_lU_l^\top)
 =\operatorname{tr}\!\left[R_l(P_l^\top A_lP_l)R_l^\top\right].
\]
내부 계산은 요청 SUM, 보고는 $J/B$로 한다.
Native owner별 reduction 뒤에 합산하며 microbatch 크기로 정규화를 바꾸지 않는다.
비용은 현재 후보의 actual key에서 만든 ideal FP64 $U$의 에너지다.
Terminal에서는 FP32 effective $U$의 비용과 차이를 따로 기록한다.
$E=\norm{R-UK}_F^2$는 진단에만 쓰고 pulse나 추가 보존 목적을 넣지 않는다.

Task와 $Q$ 모두
$u\to R\to U_{\rm lower}\to K_{\rm upper}\to P_{\rm upper}
\to U_{\rm upper}$를 미분한다.
대칭 $A$에서 하층 $j$의 비용 gradient에는
\[
 \nabla_{R_j}\sum_l Q_l
 =2R_j(P_j^\top A_jP_j)
 +\sum_{l>j}\left(\frac{\partial K_l}{\partial R_j}\right)^*
                 \frac{\partial Q_l}{\partial K_l}
\]
가 포함된다. Candidate $K/P$를 detach하지 않는다.
$G$는 prox에서 한 번만 처리하고 smooth gradient에 중복해서 더하지 않는다.

\section{한 번 정하는 $\lambda_Q$: native 기준점의 방사 방향 균형}
리뷰의 clamp 기준 제안은 \textbf{한 기준점에서 actual 목적의 방사 방향을
맞추는 가격}으로 한정해 채택한다. Native clamp와의 등가성은 주장하지 않는다.

\begin{enumerate}
 \item B1의 같은 entry와 native 문장으로 native anchor 한 층
 (현재 L8)의 compute-z를 요청마다 독립 실행한다.
 원래 Adam, norm, clamp, early exit, 최대25평가를 쓴다.
 반환된 실제 endpoint를 $R^\star$로 두고 나머지 후보층은 0으로 둔다.
 내부 endpoint를 경계까지 확대하지 않는다. Reference는 commit하지 않고 $H$도 바꾸지 않는다.
 \item 이 $R^\star$를 같은 actual MEMIT candidate로 평가한다.
 Native subject-hook gradient를 대신 쓰지 않고
 actual $F_{\rm task}$, $G$, $Q$의 방사 방향량을 측정한다.
 \item $u^\star$ 방향에 대해
 \[
 t=-\langle\nabla F_{\rm task}(u^\star),u^\star\rangle,\quad
 n=G(u^\star),\quad d=\langle\nabla Q(u^\star),u^\star\rangle
 \]
 로 두고
 \[
 \boxed{\lambda_Q=(t-n)/d}
 \]
 를 계산한다. 비영 층이 하나이므로 해당 층의 key는 이 확대 방향에
 의존하지 않으며 $d=2Q(u^\star)$도 검산한다.
\end{enumerate}
$t,n,d,\lambda_Q$가 모두 유한하고 $d>0$, $t>n$, $\lambda_Q>0$이어야 한다.
실패하면 invalid energy, nonpositive price, nonfinite/overflow를 구분해 보고한다.
Abs, epsilon 분모, clipping, 기본값으로 main fit을 시작하지 않는다.
Native 계약상 모든 요청이 실제 no-op인 경우에만 무편집 batch와
native history append를 허용하고, 첫 nontrivial batch에서 같은 규칙을 적용한다.
미수렴이나 오류로 반환된 0을 no-op으로 처리하지 않는다.

이 가격은 전체를 합친 방사 방향 미분만 0으로 만든다.
접선, 다른 층, 요청별 KKT나 의미적 강도를 맞추지는 않는다.
요청별 radial residual, 다른 층의 gradient와 상한 접촉률도 기록한다.
한 번 정한 가격은 모든 batch에 고정하고 PS/NS/ACC로 선택하거나 다시 맞추지 않는다.
\textbf{Native 기준점 fit은 추가 계산}이다. 3점 sweep은 아니지만 무료도 아니므로,
기준점 생성과 actual gradient 계산 시간을 main과 구분해 총비용에 포함한다.

\section{Solver: 상한을 둔 nonmonotone proximal BB1}
기존 EfficiencyAdam을 대신하여 $g=\nabla F(u)$에 대해
\[
 v_{lr}=u_{lr}-\eta g_{lr},\qquad
 u^+_{lr}=
 \frac{\min\{c_{lr},(\norm{v_{lr}}-\eta\beta_{lr})_+\}}
      {\norm{v_{lr}}}\,v_{lr}
\]
를 사용한다. $v=0$이면 0으로 정의한다.
$u$ 좌표 step을 $\eta_u=\eta$라 두면 문턱값은 $\eta_u\lambda_n/a_{lr}$이다.
같은 갱신을 $R$ 좌표로 옮길 때는 block별
$\eta_{R,lr}=\eta_u a_{lr}^2$이므로 서로 다른 좌표에 같은 step을 대입하지 않는다.
정확한 0이 영구 비활성화를 뜻하지 않는다.
다음 full gradient에서 모든 block을 평가하여 재진입을 허용한다.

초기 step은 $\eta_0=\min c_{lr}/\max_{lr}\norm{g_{lr}(0)}$로 정한다.
초기 $u=0$에서 모든 $g=0$이면 합성 목적의 mapping도 0이므로
$\eta_0,\tau$는 해당 없음으로 기록하고 정상성 종료한다.
그 외에는 $s=u_k-u_{k-1}$, $y=g_k-g_{k-1}$에서
유한한 $s^\top y>0$일 때 BB1 $\eta=s^\top s/(s^\top y)$를 사용한다.
아니면 직전 accepted step을 재사용한다.
Step 범위는 $[10^{-6}\eta_0,10^6\eta_0]$로 제한한다.
$y$에는 norm subgradient를 넣지 않는다.

초기 window를 $J(u_0)$ 하나로 시작한다.
직전 accepted $J$ 최대5개의 최댓값을 $J_{\rm ref}$로 두고
\[
 J(u^+)\le J_{\rm ref}
 -\frac{10^{-4}}{2\eta}\norm{u^+-u}_F^2
\]
이면 채택한다. 아니면 $\eta$를 절반으로 줄여 다시 평가한다.
Rejected 후보도 fresh $K/P$와 actual forward가 필요하지만 backward는 하지 않는다.
하한에서 계속 불충족하거나 같은 trial에 정체되면
\texttt{NUMERICAL\_STALL}로 구분한다.

\begin{center}
\begin{tabular}{lr}
\toprule
Main solver의 상한 & 값\\
\midrule
Accepted updates & 24\\
Logical candidate evaluations (초기와 rejected 포함) & 50\\
Full gradient evaluations (초기와 terminal 포함) & 25\\
\bottomrule
\end{tabular}
\end{center}
초기 gradient 한 번과 accepted update당 한 번으로 최대25회다.
Terminal은 마지막 accepted gradient를 재사용하며 별도 26번째 backward를 하지 않는다.
Gradient를 계산하기 위한 replay forward는 물리 호출로 따로 센다.
Calibration, qualification, observer도 별도 counter로 세되 총비용에는 포함한다.
50개 후보는 물리 forward50회라는 뜻이 아니다.
마지막 \textbf{accepted} 후보를 반환하고 마지막 rejected 후보를 commit하지 않는다.

\paragraph{정상성과 미수렴.}
고정 probe step $\tau=\eta_0$에서
\[
 {\cal G}_\tau(u)=
 \frac{u-\operatorname{prox}_{\tau G+I_{\mathcal C}}
                    (u-\tau\nabla F(u))}{\tau}
\]
를 accepted 후보와 terminal에서 기록한다.
각 block의 $\norm{{\cal G}_{\tau,lr}}/
\max(1,\norm{\nabla F_{lr}},\beta_{lr})$ 최댓값이 $10^{-3}$ 이하이면
\texttt{STATIONARY\_TOL}, 상한에서 미달하면
\texttt{BUDGET\_NOT\_STATIONARY}로 보고한다.
Raw max/RMS/Frobenius, 분모와 $\tau$도 함께 기록한다.
Native all-loss$<.05$를 이 공동 composite 목적의 종료 조건으로 쓰지 않는다.

미수렴이어도 유한값, feasibility, 후보/commit 정합성을 통과한 고정예산 후보는
기존 사용자 방침에 따라 반환할 수 있다. Calibration, NaN, solve, identity 오류는 별개다.
FP32 cast/add의 backward는 smooth surrogate이므로 production 잔차를
양자화된 함수의 엄밀한 KKT 증명이라고 부르지 않는다.
작은 잔차도 전역 최적해나 편집 성공을 뜻하지 않는다.

\section{실현률, history와 누적 손상에 관한 한계}
대칭 SPD $A$에서 $S=K^\top A^{-1}K$이면
\[
 UK=R\,S(I+S)^{-1},\qquad
 Q=\operatorname{tr}[R\,S(I+S)^{-2}R^\top]
 \le \tfrac14\norm{R}_F^2 .
\]
Requested share와 realized share가 같아야 하는 것은 아니다.
후보가 평가한 실제 작용과 commit의 작용이 같아야 한다.
한 층의 부족분을 다른 층의 tracking으로 강제 보충하지 않는다.
위 batch 행렬 성질을 요청별/context별 ratio 상한으로 읽지 않는다.

\textbf{History가 늘면 $Q$가 반드시 줄어든다는 예측은 채택하지 않는다.}
단일 key와 고정 $R$에서도
$Q=\norm{R}^2\kappa/(1+\kappa)^2$이므로
$\kappa:7\to1$이면 $Q/\norm{R}^2:.109375\to.25$로 증가한다.
Batch마다 최적화한 $R/K$도 달라진다.
따라서 $Q$ 비증가나 NS 손상의 선형화를 보장하지 않는다.
$Q_{C_0}$, $Q_H$, 총 $Q$, $Q/B$, 실현량, ACC, cohort retention을 나눠 측정한다.
낮은 $Q$만으로 의미적 보존을 주장하지 않는다.

\section{Commit, 계산량과 모델ㆍbatch 일반성}
마지막 accepted actual forward가 사용한 같은 FP32 weight payload를 copy해 commit한다.
새 key로 다시 solve하거나 같은 update를 두 번 더하지 않는다.
모든 층 write 후 native rewrite mean key로 history를 한 번만 append한다.
마지막 후보의 causal key를 재사용할 경우 final model과의 parity를 먼저 검증한다.
후보 평가 중 persistent $W/H$를 바꾸지 않고,
durable checkpoint는 기존 방침대로 만들지 않는다.

변하지 않는 prefix와 batch-entry $A$ factor를 재사용한다.
Actual causal builder를 native loss readout까지 이어서 별도 subject-only forward를 없앤다.
Whole-batch barrier를 보존하는 stage replay/offload, compact $Q$,
필요한 위치만의 full-vocabulary head와
dense $dW$를 만들지 않는 $R/P$ 직접 VJP를 사용한다.
Rejected 후보는 forward-only다.
이 절약은 loss, K/P gradient, token/문장, 논리 batch를 바꾸지 않는다.
상층 key의 entry 고정이나 detach를 시간 절약에 섞지 않는다.

Adapter는 층 순서, 차원, write 위치, native context/lookup/readout/reduction을 제공한다.
L4--L8, hidden4096, BS100, 일정 context 수를 알고리즘에 고정하지 않는다.
물리 microbatch는 실행 단위이며 writer의 논리 batch를 나누지 않는다.
임의 batch shape를 지원한다는 사실이 batch를 바꿔도 같은 해가 나온다는 뜻은 아니다.
속도 향상과 2k ETA는 미측정이며 V14의 후보 시간에서 바로 외삽하지 않는다.

\section{B1과 2k의 검증 및 판정}
기존 CPU full-causal reference 12개와 이번 calibration/prox 수학 검증12개는 통과했다.
Production BB controller, 실제 모델에서 양의 가격 존재, FP32 정합,
성능과 runtime은 아직 검증하지 않았다.
구현 후 2--4개 요청에서 dense/streamed 값ㆍfull gradient,
유한차분, raw native solve, zero/reentry, controller 상한,
같은 payload의 commit, history-once를 먼저 검증한다.
대칭 SPD는 에너지 정리의 가정이며 production의 raw native $A$를
몰래 대칭화하거나 jitter를 넣지 않는다. 같은 $A$의 solve/VJP와 수치 오차를 검증한다.

주 실험은 같은 ordered2,000요청, BS100$\times$20, fresh $W_0/H_0$의 단일 MEMIT arm이다.
B1의 PS95\%, NS 손실0.4pp, CD보다 낮은 $Q$는 \textbf{비교 목표}다.
서로 다른 방법의 최선 수치를 합친 엄격한 진급 조건으로 만들지 않는다.
평가 성능으로 계수나 후보를 고르지 않고 B1을 재fit하지 않는다.
작은 technical 검증과 계수 정의의 성립을 과학적 성능의 성공과 구분한다.

W0, 각 cohort의 pre/post, W5/W10/W15/W20의 누적 및 고정 cohort를 평가한다.
RS/PS/NS의 분자ㆍ분모, TF token/prompt/strict ACC,
new/true NLL, W0-correct retention, own-write/later-write 손상을 보고한다.
W10 이후 NS slope는 PS 및 strict ACC와 함께 읽는다.
Active layer, 재진입, 상한 접촉, requested/realized share,
$\lambda_Q$ receipt, 정상성 잔차와 미수렴률,
물리 F/Bㆍsolveㆍreplayㆍ시간ㆍmemory도 저장한다.
MEMIT-BLUE, MEMIT-H, CD와의 역사 비교에는 history, runtime, 평가 조건의 차이를 명시한다.

\clearpage
\part*{역사 기록: CD-Q / CD-C와 SH4 2,000-edit 실행}
\paragraph{권한과 범위.}
사용자의 ``구체화 진행한 이후 GH에게 전달하자. sh4에게 task 진행하라고 해.
2000edit 진행하자''를 실행 근거로 한다. GH가 SH4에 두 arm을 각각 BS100 $\times$20으로
배정한다. B1은 같은 연속 실행의 첫 batch이며 성능에 따른 진급 심사가 아니다.
현재 Server3 V14 실행은 별도다. 이 문서는 기존 편집기의 같은 파일을 갱신하며,
아래 V14 및 V13 본문은 역사 기록으로 보존한다. 실행 계약은
\texttt{../2026-10-05-jlz-cd-cumulative-allocation/}의
\texttt{method-contract.json}, \texttt{implementation-contract.json},
\texttt{experiment-2k.json}이다. Production 구현과 GPU 성능은 SH4의 후속 작업이다.

\section{교정 대상과 유지하는 backbone}
V13 CD는 현재 subject의 직접 변위는 거의 정확히 실현했으나 누적 locality와
정답 ACC의 손실이 컸다. 다음 설계는 이 현상을 history 부재나 수치 조건수 폭주로
단정하지 않는다. 목표는 정확히 실현할 local 목표를 정할 때 실제 write 부담과
기존 변화와의 간섭을 함께 고려하는 것이다. 모든 후보 층을 유지하며 특정 층에
배분이 집중되는 결과도 허용한다. Native 문장, subject lookup, NLL reduction,
current$\Vert$entry KL, V13 norm 계수, EfficiencyAdam과 요청별 공유 예산 .75를 유지한다.
새 reference, generation 문장, replay, L8 completion이나 tracking은 넣지 않는다.
CD writer의 compatible exact / retained range-LS 계약과 actual 상층 key 재측정은 유지한다.

\section{고정 entry CD 연산자의 축약}
Batch의 실제 요청 수를 $B$, 층 $l$의 native injection row 수를 $m_l$이라 한다.
$D_l\in\mathbb R^{d_l\times B}$는 requested residual, $E_l\in\mathbb R^{B\times m_l}$는
명시적인 owner 행렬이다. 기존 CD row 가중치를
$\Omega_{jj}=1/(B n_{\mathrm{owner}(j)})$로 두되 NLL의 native reduction과 구분한다.
현재 entry key $K_l$, covariance와 history로
\[
 A_l=\lambda_{C,l}C_{0,l}+H_l=L_lL_l^\top,\qquad
 X_l=L_l^{-1}K_l\Omega_l^{1/2}
     =P_l\Sigma_lV_l^\top
\]
를 계산한다. 식의 SVD는 기존 CD의 retained 부분이며 cutoff를 새로 강화하지 않는다.
Inverse는 triangular solve와 retained singular value 나눗셈으로 구현한다.
\[
 \mathcal B_l=E_l\Omega_l^{1/2}V_l\Sigma_l^{-1}P_l^\top L_l^{-1},\quad
 \widehat U_l(D_l)=D_l\mathcal B_l,\quad
 S_l=\mathcal B_lK_l,
\]
\[
 F_l=\mathcal B_lA_l\mathcal B_l^\top,\qquad
 \widehat Q_l(D_l)=\operatorname{tr}(D_lF_lD_l^\top).
\]
$F_l$은 $B\times B$이고 off-diagonal을 유지한다. 이를 대각선으로 치환하면
요청 사이 항을 버린 다른 근사 method다. $K,A,\mathcal B,S,F$는 batch fit 동안 고정한다.

Fit의 subject 주입은 $\widehat Y_l=D_lS_l$이다. Compatible target에서 값은
$D_lE_l$과 같고, incompatible target에서는 CD가 실제로 만드는 retained response다.
현재 $D$에서 값이 같다는 사실만으로 $S$의 Jacobian을 $E$로 바꾸지 않는다.
모든 owner 방향에서 $S=E$가 검증된 경우만 direct-D fast path를 허용한다.
이 계약은 실현되지 않는 방향이 virtual loss만 개선하는 허점을 막는다.
Actual commit의 상층 key는 entry key와 다르므로 이를 완전한 actual 경로라고 부르지 않는다.

\section{누적 간섭 비용과 해석}
최초 모델의 selected weight를 $W_{\mathrm{initial},l}$로 보존하고,
실제 FP32 committed weight로
\[
 \Delta_l=W_{\mathrm{entry},l}-W_{\mathrm{initial},l},\qquad
 J_l=\Delta_l C_{0,l}\mathcal B_l^\top,\qquad
 \widehat c_l(D_l)=\langle D_l,J_l\rangle_F
\]
를 만든다. 이상적 update의 누적합을 실제 weight 차이로 대체하지 않는다.
두 arm의 비용은
\[
 \boxed{\Pi_l^{(\alpha)}(D_l)=\widehat Q_l(D_l)
          +2\alpha\lambda_{C,l}|\widehat c_l(D_l)|},\qquad
 \alpha=0\ (\mathrm{CD\!\!\;Q}),\quad \alpha=1\ (\mathrm{CD\!\!\;C}).
\]
절댓값은 층마다 취한 뒤 합한다. 부호 있는 교차항을 직접 더하면 과거 편집을
취소하는 방향이 음의 비용을 얻을 수 있다. 절댓값은 이 보상을 제거하지만,
정당한 후속 수정이나 유익한 상쇄도 벌할 수 있다.

고정 $C_0$에서
\[
 g_l=\norm{(\Delta_l+\widehat U_l)C_{0,l}^{1/2}}_F^2
      -\norm{\Delta_l C_{0,l}^{1/2}}_F^2
\]
이면 triangle inequality로
\[
 \lambda_{C,l}|g_l|
 +\operatorname{tr}(\widehat U_lH_l\widehat U_l^\top)
 \le \Pi_l^{(1)}
\]
이다. 이는 covariance 변위의 proxy bound이지 의미적 손상이나 ACC 보장이 아니다.
$\Delta$에는 성공한 편집도 포함되고 한 층 안의 서로 다른 작용은 상쇄될 수 있다.
B1은 $\Delta=0$이므로 두 arm의 목적과 후보 경로가 같아야 한다.

\section{공동 최적화, 계수 및 종료 계약}
V13의 $D_{lr}=a_{lr}u_{lr}$와 $\sum_l\norm{u_{lr}}_2\le .75$를 유지한다.
\[
 J_r^{\mathrm{native}}=\mathrm{NLL}_r(\widehat Y)
 +\lambda_{\mathrm{KL}}\mathrm{KL}_r(\mathrm{current}\Vert\mathrm{entry})
 +\frac{.5}{a_r^\star}\sum_l\norm{u_{lr}}_2,
\]
\[
 J_{\mathrm{sum}}=\sum_rJ_r^{\mathrm{native}}
       +\lambda_{\mathrm{alloc}}\sum_l\Pi_l^{(\alpha)}.
\]
Backward는 sum, 보고는 $J_{\mathrm{sum}}/B$로 하여 optimizer scale을 보존한다.
고정 geometry의 추가 gradient는
\[
 \partial_{D_l}\Pi_l^{(\alpha)}
 =2D_lF_l+2\alpha\lambda_{C,l}\operatorname{sign}(\widehat c_l)J_l,
 \quad\operatorname{sign}(0)=0.
\]
$u$로 전달할 때 열별 anchor를 곱한다. Native norm gradient는 한 번만 더한다.
모든 요청의 같은 candidate에서 native 및 비용 gradient를 합한 뒤 갱신한다.
Projected response도 요청 사이 의존성을 만들 수 있으므로 request freeze를 남겨 두지 않는다.
모든 $J_r^{\mathrm{native}}<.05$이면 공통 종료, 최대25평가/24갱신,
마지막 평가 후보 채택 및 terminal backward 없음으로 고정한다.
이 종료는 risk 목적의 stationary point 증명이 아니다. 기존 V13 독립 종료와의
차이를 역사 비교에 명시하며 두 신규 arm에는 같은 제어를 적용한다.

\paragraph{배분 계수의 사전 규칙.}
B1의 첫 비영 후보(통상 c1)에서 $u$ 좌표의 전체 native norm gradient $g_n$과
$\widehat Q$ gradient $g_Q$를 계산하여
\[
 \lambda_{\mathrm{alloc}}=\norm{g_n}_F/\norm{g_Q}_F
\]
로 한 번 정하고 두 arm 및 전체20 batch에서 고정한다. Ratio1은 비용 단위를 정하는
명시적 관례이며 최적 강도의 보장이 아니다. Native norm .5는 바꾸지 않는다.
c0에서 $D=0$, $\Delta=0$이라 비용 gradient는0이므로 첫 update에 별도 warm-up은 없다.
추가 fit이나 평가 성능 sweep도 없다. 모두0이면 첫 비영 후보까지 pending하며,
전 batch가 무편집일 때만 $\Delta=0$ 검증 후 다음 batch로 pending을 넘긴다.
비영 실현에 $g_Q=0$ 또는 비유한 값이 나오면 기술 오류다. 임의 epsilon, 계수 clip,
default 값으로 숨기지 않는다. Calibration 값과 source/candidate/geometry hash를 lock한다.

\section{실행 비용과 actual commit}
기존 initial capture를 재사용하므로 geometry 준비에 새 모델 forward는 필요 없다.
Batch마다 entry QR--SVD를 층 수만큼 추가하고, 변하지 않는 A factor를 commit과 공유한다.
첫 층 key의 동일성을 확인하면 해당 SVD도 재사용할 수 있어 현재5층은
총5회에서9회 SVD가 된다. 상층 actual key 재측정과 solve는 유지한다.
반복25회에는 actual builder 또는 solver backward를 추가하지 않는다.
현재 B100, 5층의 $D F$는 후보당 약0.410 GFLOPs지만 entry의 $C_0\mathcal B^\top$
준비는 약205.5 GFLOPs이므로 실제 FP64 시간과 전송량을 측정해야 한다.
Selected $W_{\mathrm{initial}}$의 FP32 CPU 저장은 현재 profile에서 약1.094 GiB다.
모든 factor를 보관하면 추가 CPU 메모리가 크므로 한 층씩 전송하고 memory cap을 지킨다.
Batch size, 문장 수, layer 수와 차원은 adapter/owner metadata에서 받는다.
Logical batch를 physical microbatch 때문에 분할하거나 rank cutoff를 바꾸지 않는다.

Commit은 V13 CD의 target $D_lE_l$를 actual fresh key에서 푼다.
예측 $\widehat Q,\widehat c,\Pi$와 실제 effective update의 값을 따로 기록한다.
Gap 분해는 fit에서 실제 주입한 $\widehat Y$에 맞춰
\[
 h^{v,\mathrm{pre}}=h^{v,\mathrm{post}}-\widehat Y,\quad
 b^{\mathrm{inh}}=h^{a,\mathrm{pre}}-h^{v,\mathrm{pre}},\quad
 e^{\mathrm{fit}}=U_{\mathrm{eff}}k^a-\widehat Y
\]
로 정의한다. $U_{\mathrm{eff}}k^a-D_{\mathrm{owner}}$는 requested realization 오차로
별도 기록한다. 단순 항등식 검산만으로 이 두 기준의 혼동을 탐지할 수 없다.

\section{SH4 실험 설계와 판정 범위}
두 arm을 각각 fresh W0/H0에서 고정 CounterFact prefix2000, BS100 $\times$20으로 실행한다.
Model/tokenizer/C0/native prompt와 입력 순서는 V13 CD 자산과 동일성을 확인한다.
FP32 model과 FP64 geometry를 유지한다. 자원은1 GPU/arm, task동시최대2 GPU이며
최신 Server4 cap 안에서 순차 실행도 허용한다. 기존 무관 job을 취소하지 않는다.
작은 CPU reference와2--4요청 GPU 정합 검증 후 성능 심사 없이20batch를 진행한다.
B1은 두 arm의 동일성 확인과 사실 보고 시점이며 별도 reset 실험이 아니다.

W0에는 전체2000을 평가하고, 각 batch는 current cohort pre/post를 평가한다.
W5/W10/W15/W20에는 누적 prefix, 동일 first500과 birth-cohort retention을 집계한다.
동일 endpoint의 문항은 중복 실행하지 않는다. RS/PS/NS 외에 TF token/prompt/strict ACC,
new/true NLL, 기존 정답 lost/gained/retained를 남긴다. 현재 benchmark W20 분모는
R2000/P4000/N20000이며 adapter로 실제 수를 검증한다.
예측--실제 비용 차이, requested/realized share, 총변위, budget, 계수 lock,
exclusive 시간과 peak memory를 함께 기록한다. 평가 문장을 fit/계수/정지에 쓰지 않는다.

주 비교는 같은 runtime의 CD-Q 대 CD-C다. 기존 V13 CD/MD와 native baseline은
동일성 표를 갖춘 historical reference이며 새 baseline 제출은 이 지시에 포함하지 않는다.
NS만 오르고 PS/P strict가 크게 낮아지면 단순 강도--보존 이동일 수 있다.
배분 개선 주장은 품질 유지, 고정 cohort 누적 손실 감소, 실제 비용과 배분 대응이
함께 관측될 때 검토한다. SH4는 수치와 사실만 보고하고 GH가 해석한다.

\clearpage
\part*{역사 기록: 기존 V14 실행 교정과 V13 설계}
\paragraph{V14 기록 당시의 범위.}
아래 기록은 V14 B1 및 V13 2k 완료 결과를 반영한 V14 실행 교정 R2이다.
기존 열린 편집기와 파일 경로를 유지하며 V13 원 설계는 뒤의 역사 부록에 보존한다.
기준 method는 V14의 native ridge와 공동 requested-residual 최적화다.
기존 .75는 실행 동등성 검증의 기준이며, 최신 사용자 선택으로 다음 실험에는
공유1.5와 층별.75를 적용한다. 아래 Server3 2k 실행 절이 이 예산을 명시적으로 override한다.
R2는 중복 계산과 불필요한 adjoint를 줄이는 실행 설계이며, production 구현이나
새 GPU 실험을 수행했다는 뜻이 아니다. Companion은
\texttt{../2026-10-05-jlz-v14-native-writer-aware/}의 method/runtime 계약과
\texttt{execution-optimization-r2.json}이다. 과거 실행의 source lock은 변경하지 않는다.

\section{검토 결과와 유지할 연구 대상}
실제 V14 B1 source는 \texttt{2ab04d0b}, 완료 결과의 게시 기준은
\texttt{2fdd23f3}이다. V14는 같은 환경의 V12 방식 tracking 대조군에 비해
PS 88.5\%에서 90.5\%, P strict ACC 56.5\%에서 63.5\%, NS 87.8\%에서
88.2\%로 개선됐다. 반면 canonical 정규화 직접 변위의 합은 .7532에서 .4869로
작아졌다. 이를 native writer 자체의 실현률 향상으로 해석하지 않는다.
동일 direct ridge 대조군 대비 PS 73.5\%에서 90.5\%로의 개선은 실제 writer 응답을
고려하는 목표 생성의 근거다. 이 B1 관측만으로 순차 locality를 보장하지 않는다.

V13은 현재 context의 국소 목표를 정확히 실현했으나 2k에서 누적 보존이 악화됐다.
V14는 native soft ridge의 fit--보존 절충을 유지한다. 연구 claim은 모든 후보 층의
local 요청량을 jointly 계획하고, actual writer 응답에 따라 배분하는 것이다.
최소 semantic damage, 최소 writer 에너지 또는 globally optimal allocation을 주장하지 않는다.
모든 eligible 층을 고려하며 특정 층으로 집중되는 결과 자체를 실패로 판정하지 않는다.

\section{수학적 method: R2에서 변경하지 않는 계약}
Batch entry의 $W^0,H,C_0$, native teacher, canonical anchor $a_{lr}$와
native anchor $a_r^\star$를 고정한다. 모든 eligible 층에
\[
 R_{lr}=a_{lr}u_{lr},\qquad
 \sum_l\norm{u_{lr}}_2\le c,\qquad c=.75
\]
를 둔다. 각 candidate는 같은 entry에서 시작하며 이전 candidate의 write/history를
누적하지 않는다. 실제 하층의 all-token write를 거친 현재 native key를 모아
rewrite-only native nested mean $K_l$을 구성한다.
\[
 A_l=\lambda_C C_{0,l}+H_l,\quad
 P_l=(A_l+K_lK_l^\top)^{-1}K_l,\quad U_l=R_lP_l^\top.
\]
현재 profile의 $\lambda_C$는15000이며 다른 profile은 native 설정에서 받는다.
행렬 inverse는 solve로 구현한다. Native 저장 A와 dtype/order를 유지하며 임의의
대칭화, jitter, exact writer 또는 사영 writer로 치환하지 않는다.
실제 forward는
\[
 W_l^{eff}=W_l^0+\operatorname{FP32}(R_{l,64}P_l^\top),\quad
 v_{lrc}=\operatorname{linear}(k^a_{lrc},W_l^{eff})
          -\operatorname{linear}(k^a_{lrc},W_l^0)
\]
이다. 이 context별 actual subject 작용을 별도의 native subject-only 경로에 주입한다.
Native rewrite/KL 문장, lookup, target-token reduction 및 full-vocabulary readout은 유지한다.
\[
 J_{\mathrm{sum}}(u)=\sum_r\left[
 \mathrm{NLL}_r(v(u))+.0625\,\mathrm{KL}_r(\mathrm{current}\Vert\mathrm{entry})
 +\frac{.5}{a_r^\star}\sum_l\norm{u_{lr}}_2\right].
\]
Gradient는 모든 요청의 합이며 보고는 요청 평균이다. Native mean, $P$, actual upper
keys 및 하층 경로 전체를 미분한다. Norm gradient는 해석적으로 한 번만 더한다.
EfficiencyAdam의 요청별 상태, FP64 projection/FP32 저장, zero-layer reentry를 유지한다.
모든 $J_r<.05$일 때만 공통 early exit를 허용하고 최대25평가/24갱신,
마지막으로 평가한 유한 candidate의 weight payload를 commit한다. 채택 terminal에서는
backward를 하지 않는다. 마지막 history는 모든 층 write 후 rewrite 평균 key로 정확히 한 번
append하며 KL이나 새로운 reference 문장은 넣지 않는다.

\paragraph{예산 확대의 채택과 해석.}
최신 사용자 지시로 다음 실험에는 공유1.5와 층별.75를 채택한다. 이는 feasible set과
norm/stop 관계를 바꾸는 과학적 변경이며 실행 최적화와 분리해 기록한다.
두 층의 norm 배분 형태를 허용할 뿐 BLUE의 순차
local-z, 상태별 anchor, optimizer, history까지 정확히 포함하지 않는다.
낮은 B1 $Q$는 확대 후 locality의 안전 여유가 아니다. Native-only 조건에 따라
\texttt{generation\_prompts} 추가도 이 교정에 포함하지 않는다.

\section{실측 계산량과 최적화 순서}
같은 Server4 환경에서 V12 계열 planner fit은385.585초, V14 fit은2038.636초로
5.287배다. V14 전체 B1 GPU allocation은2439초였다. 보존 비용 $Q$와 GPU 계산량을
구분한다. 다음 값은 서로 겹치지 않는 원 candidate 타이머의 합이다.
\begin{center}
\begin{tabular}{lr}
\toprule
구간 & 초\\\midrule
Actual candidate builder & 107.037\\
Native 최초 정지 평가 & 234.271\\
Native replay forward/backward & 962.572\\
Builder reverse & 722.109\\
Optimizer & 6.942\\
기타 & 약5.705\\\bottomrule
\end{tabular}
\end{center}
중복 forward 제거 대상은 비terminal24회의 약224.8초, 기존 fit의 약11\%다.
이 한 변경으로 5배 가속을 예상하지 않는다. 주 병목인 native F/B 및 builder reverse에는
microbatch와 cached adjoint 개선이 별도로 필요하다. A factor/prefix 재사용,
selected-position head, microbatch별 오른쪽 padding 제거, materialized weight 공유,
dense $dW$를 만들지 않는 direct $R/P$ VJP는 이미 적용된 것으로 신규 절감에 세지 않는다.

\subsection{R1: 자원 한도에 따른 physical microbatch}
논리 batch $B_t$와 실행 row/token/activation-byte cap을 분리한다. 먼저 원 global-row
순서를 유지하며 qualified physical grouping을 선택한다. 2/4/8행은 현재 profile의
측정 후보일 뿐 method 상수가 아니다. 큰 group은 padding 토큰을 늘릴 수도 있으므로
유효/패딩 토큰, 시간, peak GPU/CPU/pinned memory를 함께 기록한다.
K의 native 평균과 ridge solve는 항상 전체 logical B를 사용한다.
OOM은 같은 candidate의 임시 상태와 RNG를 복원한 뒤 작은 physical cap으로 재시도하며
논리 B, 문장, dtype, optimizer step이나 history를 바꾸지 않는다.

Length sorting은 최초 구현에 넣지 않는다. 현재 capture가 encounter order에 의존하므로
임의 정렬은 anchor/key/owner를 바꿀 수 있다. 추후 explicit global-row scatter와 native
reduction 순서 parity가 확인될 때만 별도 적용한다.

\subsection{R2: 이미 계산한 ridge 해의 정확한 adjoint}
$B_l=A_l+K_lK_l^\top$, $P_l=B_l^{-1}K_l$라 두고, 모든 row와 상층 경로의
$G_l=\partial L/\partial P_l$을 모은 뒤 한 번
\[
 Z_l=B_l^{-\top}G_l,\qquad
 \boxed{\frac{\partial L}{\partial K_l}
 =Z_l-Z_l(P_l^\top K_l)-P_l(Z_l^\top K_l)}
\]
를 계산한다. 이 식은 A가 고정된 일반 비대칭 저장 행렬에도 유효하며, LU backend에서는
실제 transpose solve를 사용한다. Explicit inverse는 만들지 않는다.
Forward의 $K,P,Y=A^{-1}K$ 및 $I+K^\top Y$ factor를 캐시해 reverse에서 primal
ridge를 다시 구성하지 않는다. A inverse-apply backend가 검증된 경우에는
\[
 Z_l=\operatorname{solve}\bigl(A_l^\top,\ G_l-K_l(P_l^\top G_l)\bigr)
\]
로 entry의 A factor를 재사용한다. A가 singular이거나 이 경로의 residual이 기준을
통과하지 못하면 원 native $B_l$ factor의 transpose solve를 사용한다.
후보마다 큰 $B_l$를 새로 factor하는 것은 이 최적화의 기본 경로가 아니다.
새 transpose solve 자체와 upper K gradient는 남는다.
Native FP32 nested mean과 FP64 cast의 adjoint 순서도 유지한다.

Cache에는 entry W/H, row/context identity, layer, dtype/backend, K identity와
candidate의 하층 update dependency를 결속한다. A factor는 entry 내에서만 재사용하고,
첫 write층의 K/P만 인과 독립성 검증 후 candidate 간 공유할 수 있다.
상층 K/P를 이전 candidate에서 가져오는 것은 허용하지 않는다.
Primal residual 검사는 cache 생성 시 한 번 수행해 receipt와 결속하고,
새 transpose solve의 residual도 검사한다. 속도를 위해 수치 검사를 삭제하지 않는다.
Forward의 FP32 materialization을 저랭크 덧셈으로 바꾸지 않는다.

\subsection{R3: 첫 편집층의 사용되지 않는 adjoint 제거}
Adapter가 첫 eligible site의 key와 입력 boundary가 모든 최적화 변수와 독립임을
입증한 경우, 그 층의 P/K/하위 boundary adjoint는 생략한다. 단순히 계산 후 버리는
것이 아니라 custom backward의 input-gradient 필요 여부로 불필요한 GEMM과 CPU 이동을
실제로 제거한다. 상층에서 내려오는 작용을 포함한 첫 층의 R gradient는 전부 유지한다.
R=0인 층의 gradient도 유지한다. 인과 독립성이 불명확한 adapter는 기존 경로를 쓴다.
생략한 P/K adjoint는 측정된0이 아니라 \texttt{NOT\_NEEDED}로 기록한다.

\subsection{R4a: 실제 완료 요청으로 공통 정지 불가를 확인한다}
수학적인 KL 비음수성이나 norm만으로 정지를 예측하지 않는다. Native FP32 KL에는
작은 음수가 가능하고, 부분 context의 손실은 전체 요청의 손실이 아니다.
반면 같은 candidate에서 어느 요청 r의 모든 native NLL/KL 항을 원래 정밀도와
reduction으로 계산한 유한 $J_r\ge.05$는 공통 early stop이 거짓이라는 정확한 증거다.

초기 권장 구현은 작은 no-grad probe다.
\begin{enumerate}
\item 마지막 cap candidate이면 기존 full no-grad 평가 한 번 후 backward 없이 채택한다.
\item 그 외에는 현재 norm penalty 최대 요청 하나를 고른다. 동률은 입력순이며 평가
paraphrase/locality를 선택에 사용하지 않는다.
\item 그 요청의 모든 native row가 들어 있는 원래 rowgroup을 그대로 no-grad 계산한다.
다른 모양으로 재포장하거나 retokenize하지 않는다.
\item 완성된 $J_r\ge.05$이면 이번 candidate의 비종료가 확정된다. 이제 전체 native loss를
한 번 forward/backward하며 모든 row의 v-adjoint를 합친다.
\item 증거를 얻지 못하면 probe의 detached row 값을 재사용할 수 있는 범위에서 원래
full no-grad 평가를 마치고, 원래 stop/replay 규칙으로 돌아간다.
\item 전체 loss/adjoint의 유한성과 whole-B reverse가 완료된 뒤에만 optimizer를 갱신한다.
나중 row가 비유한 경우 앞서 만든 adjoint도 폐기하며 commit/history를 하지 않는다.
\end{enumerate}
Probe는 추가 optimizer candidate나 fit이 아니라 실제 forward 작업으로 기록한다.
B=1에서는 이득이 거의 없어도 같은 의미를 유지한다. Current B1의350groups 중 약4groups만
probe한다는 균일시간 근사에서는 직접 절감 대상이 약222초지만, 구현 실측 가속률은 아니다.

이 경로는 전체 loss를 두 번 계산하던 동일성 검사를 매번 수행하지 않는다.
그 변경을 숨기지 않고, 동일 profile에서 grad-enabled/no-grad native loss parity를
qualification으로 검증한다. 원래 no-grad probe가 정지 판정의 기준이다.
Parity 미확인 profile이나 수치 불일치에는 기존 두-pass 경로를 사용한다.
새 microbatch에서 stop 경계에 가까운 경우에는 backward 전에 원래 grouping의
no-grad 평가로 판정을 재확인한다. Guard는 사전 고정한 loss 수치 오차의 전파량이며,
과학적 종료 임계값 자체를 변경하지 않는다.
Probe는 group 실행 순서를 바꾸므로 결정적이며 상태를 변경하지 않는 eval forward,
또는 원 순서와 정확히 같은 group별 RNG/state replay가 필요하다.
이를 검증하지 못한 adapter는 원래 순서의 두-pass 경로를 유지한다.

\paragraph{R4b: 선택적 bounded-graph 확장.}
추후에는 원 순서대로 graph를 제한된 메모리 안에 보관하다가 완성 요청의 stop 실패가
확인되면 보관분부터 역전파할 수 있다. 메모리 상한에 닿기 전까지 증거가 없으면 graph를
버리고 detached 손실만 남긴다. 비terminal로 판명되면 버린 prefix만 먼저 재계산하여
원래 gradient 누적 순서를 유지한다. 모든 요청이 종료하면 backward는0이다.
이 경로는 별도 replay ledger 검증 후에만 사용하며 최초 구현의 필수 조건이 아니다.

\section{모델ㆍbenchmarkㆍbatch 독립적인 구현 경계}
실제 $B_t$, eligible sites, layer width, native context 수와 가중치는 adapter/RunSpec에서
받는다. Owner를 global-row modulo7처럼 복원하지 않고 명시적인 owner, role, rewrite
reduction index/group, lookup, positions와 target metadata를 사용한다.
불균등 context 수, noncontiguous write층, 0목표, 중복 key, 작은 마지막 batch를 보존한다.
빈 batch는 fit/solve/commit 없는 no-op이다. 논리 B를 작은 독립 writer들로 나누지 않는다.
현재 production의 검증된 adapter는 Llama profile이며 모든 모델 지원을 주장하지 않는다.
새 adapter는 write와 주입의 가법 대응, key 인과 의존성, mask/position 및 native loss를 검증한다.

수정 지점은 다음과 같다. 이름은 구현 예정 interface이며 구현 완료를 의미하지 않는다.
\begin{center}\small
\begin{tabular}{p{.27\linewidth}p{.64\linewidth}}
\toprule
함수/모듈 & 변경 계약\\\midrule
prepare / inputs & RowSchedule, 명시적 owner/reduction metadata, 자원 cap\\
geometry & candidate cache와 transpose-solve VJP; 같은 native A/backend\\
builder.build & primal factor/P/cache receipt 저장, actual upper key 유지\\
builder.reverse & whole-B cached VJP, 첫 site 불필요 adjoint 생략\\
physical.backward & 필요한 R/P/input gradient의 실제 연산만 수행\\
subject / optimize.fit & complete-owner no-grad probe, 공통 stop, bounded fallback\\
telemetry / collector & logical work와 probe/replay/physical work 분리\\\bottomrule
\end{tabular}\end{center}

\section{검증ㆍ적용 순서와 계산량 보고}
먼저 R1 microbatch, R2 cached VJP, R3 첫 층 pruning, R4a probe를 각각 검증한 뒤 결합한다.
선택적인 R4b는 그 이후다. CPU reference는 FP64에서 dense solve/finite difference와
cached VJP를 비교한다. SPD와 비대칭 A, 중복/zero 열, B가 feature width보다 큰 경우,
하층 R에서 상층 K/P로 이어지는 gradient, 불균등 owner rows와 tail을 포함한다.
잘못된 transpose/상층 detach를 negative control로 탐지해야 한다.

Stop scheduler는 모든 요청 성공, 정확히 threshold와 같음, 마지막 cap, 여러 group에
나뉜 owner, 작은 음수 KL, graph eviction/replay, 뒤늦은 nonfinite를 검사한다.
CPU 참조의 통과는 GPU FP32 모델의 수치 정합이나 실제 가속 검증을 대신하지 않는다.
현재 NumPy FP64 참조는15개 검사를 통과했다. Cached A transpose 식은 dense
$B^\top$ solve와 최대 $3.33\times10^{-16}$ 차이였으며, causal 두 층의 전체
gradient 유한차분 오차는 최대 $6.92\times10^{-10}$이었다.
실제 factor 재사용 구현, GPU cast/kernel parity 및 가속률은 아직 검증 전이다.

실제 모델 qualification은 기존 허용오차를 유지한다: 요청별 loss atol $10^{-5}$,
rtol $10^{-4}$; R-gradient RMS 차이는 $10^{-6}+10^{-3}$ reference RMS 이하;
native solve relative residual $10^{-8}$, FP32 budget violation $10^{-6}$ 이하.
Actual local action과 accepted payload/history 검사도 유지한다.
Zeroㆍinteriorㆍ예산 경계의 같은 후보, 여러 physical group 크기, 짧은 optimizer 전이에서
u/R, projection, 공통 stop 및 commit을 비교한다. 허용오차는 결과를 보고 넓히지 않는다.
일반 microbatch 변경의 bitwise 동일성을 주장하지 않으며 경계 판정에는 원 경로를 쓴다.

Logical candidate/update 수와 physical graph/no-grad/probe/replay F/B, valid/padded token,
solve/VJP, transfer byte, cache hit/miss, OOM retry, memory peak를 따로 기록한다.
앞의2038.6초와 동일 workload의 실측을 비교해 새 ETA를 계산하며 아직 가속 배율은
\texttt{NOT\_MEASURED}다. 원 B1 속도의2k 외삽은 fit만11.33GPUh,
전체13--14시간/arm 정도의 계획값이다. 최적화 후에도 같다고 가정하지 않는다.
추가 입력, 예산1.5, upper-key 고정, solve-gradient detach, 후보 감소는 이 실행
교정의 절감으로 세지 않는다. 최신 실행 승인의 구체적 범위는 다음 절을 따른다.

\section{사용자 승인 실행: Server3, 공유1.5, 총2000 edit}
사용자는 공유 예산 확대 실험을 Server3에 맡기고2000 edit까지 진행하도록 지시했다.
GH가 SH3에 구현ㆍ수치 검증ㆍ실행을 전달한다. 실행 정본은 companion의
\texttt{experiment-budget15-2k-server3.json}이다.
Native 입력, 손실 계수, writer/history 및 전체층 gradient는 유지하고
각 요청에 대해
\[
 \sum_l\norm{u_{lr}}_2\le1.5,\qquad \norm{u_{lr}}_2\le.75\quad\forall l
\]
를 적용한다. Norm 계수를 예산이나 층 수로 다시 나누지 않는다.
제안 step의 block norm을 $s_l=\norm{w_l}_2$라 하면 정확한 Euclidean projection은
\[
 t_l=\min\{.75,\max(s_l-\tau,0)\},\qquad
 u_l=\begin{cases}(t_l/s_l)w_l&s_l>0,\\0&s_l=0\end{cases}
\]
이다. $\sum_l\min(.75,s_l)\le1.5$이면 $\tau=0$이고, 아니면
$\sum_l t_l=1.5$가 되는 공통 $\tau\ge0$를 찾는다.
FP64 projection 후 FP32 저장에서도 두 제약을 검증한다.
단순한 전체 radial 축소와 개별 clipping을 이 projection의 대체물로 사용하지 않는다.

같은 고정 CounterFact 순서의 앞2000개를 fresh W0/H0에서 BS100으로20회 편집한다.
B1은 같은 연속 trajectory의 첫 batch이며 평가값에 따라 설정을 선택하거나 재시작하지 않는다.
기술 검증 통과 후 B20까지 이어가고, B1/W5/W10/W15/W20은 중간 보고 지점이다.
W/H는 다음 batch에 RAM으로 전달하고 기본 checkpoint 미저장 정책을 유지한다.
공통25평가/최대24갱신 및 native early stop도 유지한다.

먼저 capped projection과 R1--R4a를 단계별 검증하고, 같은 후보의 loss/full gradient,
stop, accepted payload/history를 기존 .75 및 확대 예산 양쪽에서 대조한다.
검증되지 않은 최적화 stage는 기존 동등한 계산 경로로 되돌릴 수 있지만 upper K/P
gradient, 입력, 정밀도나 과학 설정을 바꾸지는 않는다.
Server3에서1GPU 동시 상한과 현재 project 자원 한도를 확인해 제출한다.
기존 다른 task를 취소하지 않는다.

매 commit에서 누적 RS/PS/NS와 TF token/prompt/strict ACC, 새 답ㆍ원 답 NLL,
birth cohort와 retention, 층별 요청ㆍ실현 배분, shared/per-layer 제약 활성률,
$Q$, update norm, 구간별 시간과 peak memory를 기록한다.
기존 V14 .75 B1 및 baseline은 동일 cohort/분모/환경 차이를 명시해 비교하고,
별도 baseline 재실행이나 계수 sweep은 추가하지 않는다.
이번2000개 실행은 사용자 승인 상태이며 job 제출ㆍ수신ㆍ완료는 별도 receipt로 확인한다.

\clearpage
\appendix
\section*{역사 부록: V13 original method}
아래는 이전 V13 설계를 보존한 기록이다. 현재 R2의 writer/목적 계약으로 적용하지 않는다.
특히 exact/context writer는 현재 V14 native ridge와 다르다. 과거 실행 결과는 당시의
pinned source/config가 기준이며, 이 편집 문서의 변경이 실행 기록을 소급 수정하지 않는다.


\paragraph{상태와 범위.}
새 writer와 다음 실험의 설계다. 현재 실험의 실행 코드나 job을 변경하지 않는다.
CPU 참조 구현의 검증 범위와 실제 모델 qualification을 구분한다.
평균 key에서의 exact writer만으로 context 및 실제 궤적 정합까지 보장했던 이전
제안을 교정한다. 등식 제약 batch writer 자체의 신규성을 주장하지 않는다.
EMMET의 기존 등식 제약 접근과 구분되는 연구 대상은 전 층 공동 local 계획,
동적인 배분, 그 계획과 실제 국소 작용의 정합이다.

\section{고정하는 계획과 해결할 대상}
v12 planner를 유지한다. 배치 entry 모델과 anchor를 고정하고, 모든 eligible
층 $l$의 subject 위치에 $\delta_{lr}=a_{lr}u_{lr}$를 함께 주입한다.
요청별 native NLL, KL(current$\Vert$entry),
$J_r=F_r+(\lambda_n/a_r^\star)\sum_l\norm{u_{lr}}_2$와
$\sum_l\norm{u_{lr}}_2\le c$를 그대로 쓴다. 현재 profile은 $c=.75$,
$\lambda_n=.5$, $\lambda_K=.0625$, 25회 평가/24회 이하 갱신,
EfficiencyAdam과 마지막 평가 후보의 채택이다. 이는 native compute\_z 전체와
동일하다는 주장이 아니라, 직전 v12 계획의 명시적인 동결이다.
새 reference 문장, replay, pulse, writer loss를 planner에 더하지 않는다.
층의 선제 제외, 최종 비제로 층 수 강제, 예산 .47 축소도 하지 않는다.

현재 target tracking은 $R_l=z_l^v-h_l^{cur}=D_l+b_l^{inh}$를 쓴다.
따라서 계획 $D_l=[\delta_{lr}]$가 0인 요청/층도 보정 작용을 받을 수 있다.
새 direct 정책의 우변은 명확히 $D_l$이다. Writer의 등식 제약과 tracking
제거는 서로 다른 변경이며 작은 동일 계획 비교에서 각각 분리한다.

\section{제약 key와 목표를 정의한다}
각 층에 도달하면 그 branch의 실제 하층 write를 반영한 모델에서 key를 새로
수집한다. $k_{lj}\in\mathbb R^{n_l}$, owner $r(j)$와 row 역할을 기록한다.
행동 목표는 주입 공간과 writer 출력 공간의 검증된 가법 대응 아래
\[
 T_l[:,j]=\delta_{l,r(j)}\in\mathbb R^{d_l}.
\]
ModelAdapter가 이 대응, lookup, mask/position 및 native forward의 인과
의존성을 제공해야 한다. 차원만 같다는 이유로 대응을 가정하지 않는다.

\paragraph{MD: 평균 key direct 비교군.}
$K_l$의 각 열은 기존 native rewrite key의 nested 평균이다. $T_l=D_l$,
열 수는 실제 논리 batch 크기 $B_t$다. Context 평균 방식과 history 입력은
v12 MEMIT-H와 동일하다. 이 writer의 제약은 평균 key에만 적용된다.

\paragraph{CD: native-context direct 주 후보.}
$K_l$에는 native fit에서 해당 $\delta_{lr}$를 주입한 모든 row의 subject
key를 개별 열로 넣는다. 현재 profile은 rewrite 6개와 KL 1개, 요청당 7개다.
KL row에도 fit에서 같은 $\delta$를 넣으므로 그 key의 목표도 같은 $\delta$다.
KL을 새 reference로 추가하거나 그 목표를 0으로 바꾸지 않는다.
다른 adapter에서는 실제 injection-row 명세에서 열을 구성한다.
target token을 붙인 fit 입력과 native key 입력의 subject key가 같은지는
causal prefix 및 adapter parity로 검증한다. 미래 target 정보를 이용한 key를
평가 입력의 key와 같다고 가정하지 않는다.

History용 key 집합은 별개다. CD에서도 history는 기존 rewrite-only native
nested 평균을 쓴다. KL key를 history Gram에 섞지 않는다.
평가용 paraphrase/neighborhood row는 제약 및 계획에 전혀 들어가지 않는다.

\paragraph{0 목표와 순서.}
모든 원래 요청과 그 제약 열을 유지한다. $\delta_{lr}=0$인 열도 다른 요청의
write가 그 key에 작용하지 말라는 제약이다. 이를 제거하거나 요청을 deduplicate하지
않는다. 층 전체 $D_l=0$일 때만 $U_l=0$으로 solve를 생략할 수 있다.
요청별 support의 합집합은 모든 층일 수 있으므로 batch당 사용 층 감소를 약속하지 않는다.

\section{실현 우선, 그 해 중 최소 보존 에너지}
Metric은 MEMIT-H의 $A_l=\lambda_C C_{0,l}+H_{t,l}$이다.
저장 행렬의 symmetric part를 FP64로 취한 것을 정확한 metric 정의로 쓰고
antisymmetric norm을 기록한다. $A_l\succ0$를 요구하며 자동 jitter나
계수 변경은 없다. 이것은 AlphaEdit의 nullspace 제약을 유지하는 writer가 아니다.
H=0에서 A 전체에 곱한 양의 상수는 exact 해를 바꾸지 않는다. 따라서 같은
$\lambda_C$ 값을 유지해도 ridge의 fit--보존 절충 강도가 그대로 유지되는 것은 아니다.

양의 대각 가중치 $W$로 다음 사전식(lexicographic) 문제를 정의한다.
\[
 \text{첫째}\quad \min_U \norm{(UK-T)W^{1/2}}_F^2,
 \qquad
 \text{둘째}\quad \min_{U\in\arg\min\mathrm{첫째}}
                 \operatorname{tr}(UAU^\top).
\]
MD의 열 가중치는 $1/B_t$, CD는 요청 $r$의 injection row 수가 $q_r$일 때
$w_j=1/(B_tq_r)$다. 따라서 각 요청의 총 가중치는 $1/B_t$, 전체 합은 1이다.
전 열의 공통 상수는 수학적 해를 바꾸지 않지만 구현에서는 합 1 정규화를 고정한다.
이 가중치는 rank 부족 시 writer의 hidden 오차
우선순위를 정의하며 native NLL/KL의 loss 가중치와 혼용하지 않는다.
제약이 정확히 실현 가능하면 이 양의 가중치 선택은 exact 최소 에너지 해를 바꾸지 않는다.
이는 수학적 full-support 해의 성질이다. 유한정밀도에서 가중치를 바꾸면 cutoff에
대해 남는 수치 부분공간이 달라질 수 있으므로 실행에서는 위 가중치를 고정한다.

Cholesky $A=LL^\top$, $X=L^{-1}KW^{1/2}$, $Z=TW^{1/2}$를 쓰면
\[
 \boxed{U=ZX^+L^{-1}}.
\]
역행렬을 명시적으로 만들지 않는다. 이 식은 $V=UL$에 대한 최소 norm
least-squares $\min_V\norm{VX-Z}_F^2$에서 나온다.
MD가 full column rank이면 익숙한 식
\[
 U=D(K^\top A^{-1}K)^{-1}K^\top A^{-1}
\]
으로 환원된다. Singular Gram의 단순 inverse나 불안정한 $M^{-1}$ 증폭을 쓰지 않는다.

\subsection{Rank 부족도 명시적으로 정의한다}
$X=Q\Sigma V^\top$의 FP64 SVD에서
$\tau=\epsilon_{64}\max(n_l,q)\sigma_{max}$보다 큰 singular value만 유지한다.
모든 singular value가 0이면 rank=0, $X^+=0$이다.
유지한 오른쪽 basis를 $V_\tau$라 하면 실현 가능한 가중 목표는
$\widehat Z=ZV_\tau V_\tau^\top$이다. $Z=\widehat Z$이면 compatible이며,
그렇지 않으면 최소제곱의 range projection이다. 버린 작은 singular 방향은
수치 rank 계약의 일부이며 물리적으로 불가능한 방향과 구분해 기록한다.
작지만 양의 singular value를 버린 경우에는 정확히
$U=ZX_\tau^+L^{-1}$이며, whitened writer $V=UL$에
$V=VQ_\tau Q_\tau^\top$ 제약을 둔 부분공간의 사전식 최적해다.
이를 원래 full-space 문제의 무제약 최적해로 재표기하지 않는다.

같은 key/같은 목표의 중복은 compatible할 수 있다. 같은 key/서로 다른 목표는
일반적으로 incompatible하다. Rank$<q$만으로 실패시키거나 요청/층을 제거하지 않는다.
기하 상태는 compatible 또는 range-projected로 구분하며, 별도로 실제 연산
residual을 검사한다. Zero-target fast path도 별도로 표시한다.
명시한 range-LS 결과는 유한하고 연산 검증을 통과하면 실행한다. Ridge나 작은
목표로 조용히 대체하지 않는다. 큰 에너지, 낮은 PS/NS, 특정 배분은 채택 gate가 아니다.

\subsection{유한 정밀도와 실제 배포}
FP64 Cholesky 및 whitened key SVD/QR--SVD를 기본 backend로 사용한다.
$X^\top X$는 조건수를 제곱하므로 rank 판단용 기본 backend로 쓰지 않는다.
Rank와 목표 projection residual, retained singular condition number,
$\norm{UKW^{1/2}-\widehat Z}_F$를 기록한다.
이상적 solve와 별개로 native orientation과 순서대로
\[
 W_{new}=W_{old}+\operatorname{FP32}(U),\qquad
 U_{eff}=\operatorname{FP64}(W_{new})-\operatorname{FP64}(W_{old})
\]
를 적용한다(현재 FP32 profile). 실제 $U_{eff}K-T$와 모델의
$h_l^{after}-h_l^{before}-U_{eff}k_l$를 검증한다.
Algebra residual과 target approximation residual은 서로 다르다.
수치 rank가 full이더라도 실제 연산이 정확하다는 보장은 없다.
CPU 반례에서 condition 약 $10^{12}$의 retained full-rank 문제는 compatible했지만
FP64 실현 오차가 약 $9.62\times10^{-5}$였다. 따라서 실제 projected-target
항등식 검증이 실패하면 \texttt{NUMERICAL\_REALIZATION\_FAILURE}로 분리한다.
Production profile은 그 항등식에 $10^{-10}+10^{-8}\norm{Z}_F$ 허용오차를
사전 고정하고, CPU 참조의 더 엄격한 기본값과 구분한다.
실제 FP32 작용의 ideal/effective 및 native local-additivity 비교는
atol $2\times10^{-5}$, rtol $2\times10^{-4}$를 기본 qualification 값으로
고정한다. 이는 모든 원소의 절대차가 atol+rtol$\times$reference 절대값 이하인
elementwise allclose다. Effective action 대 ideal action에서는 ideal을 reference로,
실제 native hidden 차이 대 $U_{eff}K$에서는 $U_{eff}K$를 reference로 쓴다.
Target의 range 오차 자체를 이 tolerance로 gate하지 않는다.
SPD 실패, 비유한 결과, 실제 저장/복원 오류, 고정된 연산 parity 실패만 기술 중단이다.
큰 condition number나 에너지만으로 보폭/배분/목표를 자동 수정하지 않는다.

\section{어디까지 실현을 보장하는가}
CD가 compatible하면 제약에 포함된 현재 actual subject key 각각에서
$U k_{rc}=\delta_r$이다. 이는 평균만 맞추는 MD보다 강한 국소 계약이다.
모든 eligible 층에서 이 조건이 성립하고 요청의 계획 예산이 0이 아니면,
동일 entry anchor로 정의한 직접 작용 share는
\[
 \frac{\norm{U_lk_{lrc}}/a_{lr}}{\sum_j\norm{U_jk_{jrc}}/a_{jr}}
 =\frac{\norm{u_{lr}}}{\sum_j\norm{u_{jr}}}
\]
로 계획 share와 일치한다. 이는 각 실제 subject 위치의 직접 변위에 대한 조건부
보장이며, 최종 누적 hidden 변화나 task 기여의 분해는 아니다.
전체 계획이 0이면 비율은 정의되지 않으므로 share는 null, 상태는
\texttt{NO\_PLANNED\_EDIT}로 기록하고 실제 작용은 anchor 정규화 leakage로 남긴다.
같은 요청의 key 평균을 $\mu=\sum_cw_ck_c$라 하면
\[
 \sum_c w_c\norm{Uk_c-\delta}^2
 =\norm{U\mu-\delta}^2+\sum_cw_c\norm{U(k_c-\mu)}^2.
\]
즉 context 간 key 차이에도 같은 목표 작용을 요구한다. 이것은 에너지와
실현 불가능성 부담을 키울 수 있다. Feasible한 같은 actual keys/metric에서는
CD의 최소 에너지가 MD보다 작을 수 없다. 서로 다른 순차 branch의 keys가 달라진
뒤에는 이 부등식을 그대로 적용할 수 없다.

Actual/virtual 기저 차이는 별도로 남는다.
\[
 h_l^{a,+}-z_l^v
 =(h_l^{a,-}-h_l^{v,-})+(U_lk_l^a-\delta_l).
\]
Exact direct는 둘째 항을 해결한다. 모든 token에 작용한 하층 write가 만든
첫째 항까지 없애지는 않는다. Final strength 일치, locality/retention 개선,
총에너지 감소는 실험 대상이며 보장하지 않는다. .75는 계획 예산이고 실제
전체 모델 변화의 상한이 아니다. v12 배분은 writer 에너지를 직접 최소화하지 않으므로
비싼 층을 자동 회피한다고 주장하지 않는다.

\section{순차 적용, history, 범용 실행}
층 순서마다 자기 branch의 actual key를 다시 측정하고 전체 논리 batch 제약을
함께 푼 뒤 적용한다. 상층에 잔여 목표를 넘기는 divisor, terminal L8 completion,
별도 tracking 보정은 CD/MD direct에 없다.
모든 층 write가 성공한 후 모든 eligible 층의 final native mean key로
history를 한 번 append한다. $U_l=0$인 층도 history 갱신을 유지한다.
실패 시 그 batch의 W/H/RNG와 cache 상태를 원복하며 부분 commit은 게시하지 않는다.

모델, benchmark, $B_t$, row 수, 각 층의 입력/출력 차원은 adapter와 RunSpec에서
받는다. 논리 $B_t$와 fit/capture/evaluation microbatch는 독립이다. 마지막 부분
batch를 그대로 처리한다. 큰 $B$에서 제약이 feature rank를 초과하면 위 range-LS
계약을 적용하며, 작은 독립 write 여러 번으로 바꾸지 않는다. 자원 한도 내 qualified
backend가 없으면 필요한 자원을 보고하고 논리 batch를 임의 변경하지 않는다.

구현은 planner와 실제 실행 namespace를 분리한다. 변경 함수는
\texttt{capture\_native\_sites}, \texttt{build\_constraints},
\texttt{prepare\_metric}, \texttt{solve\_realized},
\texttt{apply\_sequential}, \texttt{commit\_history},
\texttt{matched\_plan\_probe}, \texttt{reduce\_telemetry}다.
기존 selected-position head와 native prefix 재사용을 유지한다. 새 실행 최적화는
별도 parity qualification을 통과한 공통 경로만 두 arm에 동일하게 사용한다.
같은 pilot entry의 $A$ factor는 자원 예산 내에서 공유할 수 있지만, 각 branch의
key/상층 activation은 다른 branch에서 재사용할 수 없다.

\paragraph{참고.}
\href{https://aclanthology.org/2024.findings-emnlp.903/}{Gupta et al., A Unified Framework for Model Editing (EMMET), 2024}.
v12 method 계약과 2026-10-05의 AlphaEdit/ridge 동일 계획 감사는 companion JSON에서 결속한다.
\end{document}
